\documentclass[11pt, a4paper]{ctexart}

% ==========================================
% 1. 基础内核与中文配置
% ==========================================
\ctexset{
	section/format = \Large\bfseries\centering,
	subsection/format = \large\bfseries\raggedright,
	subsubsection/format = \normalsize\bfseries,
	tocdepth = 2
}

\usepackage{geometry}
\geometry{left=2.5cm, right=2.5cm, top=2.5cm, bottom=2.5cm}
\usepackage{amsmath, amssymb, amsfonts, amsthm}
\usepackage{xcolor}
\usepackage[colorlinks=true, linkcolor=blue!60!black, urlcolor=blue!60!black]{hyperref}
\usepackage{enumitem}
\usepackage{fancyhdr}
\usepackage{graphicx}
\usepackage{booktabs} % [FIX]: 表格美化，提供 \toprule, \midrule, \bottomrule 命令

% ==========================================
% 2. 核心组件库
% ==========================================
\usepackage{tcolorbox}
\tcbuselibrary{skins, breakable, theorems}
\usepackage{tikz}
\usetikzlibrary{shapes, arrows.meta, positioning, shadows, calc, fit, trees}

% --- 2.1 变长编码视觉宏 ---
\newcommand{\lowfreq}[1]{%
	{\Large\bfseries\color{red!60!black} #1}%
}

% --- 2.2 思考引导盒子 ---
\newtcolorbox{thinkbox}[1]{
	enhanced,
	colback=yellow!5,
	colframe=yellow!40!black,
	title=\textbf{🤔 思考链：#1},
	fonttitle=\bfseries\sffamily,
	attach boxed title to top left={xshift=5mm, yshift=-2mm},
	boxed title style={colback=yellow!60!black, sharp corners},
	sharp corners=south,
	arc=3mm,
	drop shadow,
	breakable
}

% --- 2.3 核心策略盒子 ---
\newtcolorbox{strategybox}[1]{
	enhanced,
	colback=blue!3,
	colframe=blue!40!black,
	title=\textbf{🛠️ 工程策略：#1},
	fonttitle=\bfseries\sffamily,
	borderline west={4pt}{0pt}{blue!60!black},
	sharp corners,
	breakable
}

% --- 2.4 数学推导盒子 ---
\newtcolorbox{mathbox}[1]{
	enhanced,
	colback=gray!5,
	colframe=gray!60!black,
	title=\textbf{🧮 数学建模：#1},
	fonttitle=\small\bfseries\sffamily,
	leftrule=2mm,
	breakable
}

% --- 2.5 链接跳转宏 ---
\newcommand{\mathlink}[2]{%
	\par\noindent\hspace*{2em}\small\color{blue!60!black}
	$\hookrightarrow$ \textbf{原理深挖}：#2 (点击跳转 $\to$ \ref{#1})
}

% ==========================================
% 3. 文档正文
% ==========================================
\title{\textbf{\Huge 写作的工程学原理}\\ \Large 结构、编码与熵减的完整闭环}
\author{Structure First}
\date{\today}

\begin{document}
	
	\maketitle
	\tableofcontents
	\newpage
	
	% ==============================================================================
	% Part 0: 万物之始：为何先谈结构？ (The Primacy of Structure)
	% ==============================================================================
	\section*{0. 我先问问你：为何我要问？}
	
	在深入任何写作技巧之前，我们必须回答一个看似与写作无关的问题。
		\begin{thinkbox}{为什么我要问为什么？}
		\textbf{因为！}：动机，或者说好奇，几乎是人开始任何体系的最初/核心/全部理由
		\textbf{所以！}：接下来我会不断问问题，这些问题包含了我这本书里面每一个零件存在的理由
	\end{thinkbox}
	
	那么你翻开书，看到目录，或者打开电脑的文件夹的时候是否会想问：
	\begin{thinkbox}{为何电脑有文件夹，而不是一个巨型“文档.txt”？为何书有目录，而不是全拍一起？}
		\textbf{现象A}：如果操作系统将你硬盘里所有文件（文档、图片、程序）都平铺在一个列表里，你需要从上百万个文件中找到去年的工作报告。这无异于大海捞针，是一场认知灾难。
		
		\textbf{现象B}：如果你书的段落，目录结构啥都没有，全部粘在一起，你将要读一篇少说几万字的小作文
		
		\textbf{检索成本}：这种线性扫描的查找成本是 $O(N)$。哪怕是心爱的人的小作文，如此密密麻麻恐怕也会令人望而却步
		
		\textbf{解决方案}：文件夹系统，或者书的目录，也就是\lowfreq{树状结构 (Tree Structure)}。它将检索成本急剧降低到 $O(\log N)$。
	\end{thinkbox}
	
	写作也是一样。一篇没有清晰结构的文章，就是那个“巨型文档.txt”。再华丽的辞藻也无法拯救读者迷失在信息海洋中的痛苦。因此，\textbf{结构化是写作的第一性原理}，其唯一目的就是建立高效的\lowfreq{索引系统 (Indexing System)}。
	
	\subsection{好结构的双重质询：正交与完备}
	
	一个“好”的文件夹系统，必然是“分配合理”的。如何用工程化的语言定义“合理”？
	
	\begin{enumerate}
		\item \textbf{正交性质询 (Orthogonality)}：你的分类是否互相独立，没有重叠？
		
		\textit{反例}：在一个财务文件夹里，你同时建立了“经费”和“预算”两个子文件夹。这两个概念高度重叠（线性相关），会导致文件存放混乱，增加了检索时的不确定性。这违背了注意力原则，因为读者/用户需要思考“这个文件到底该放在哪/从哪找”。
		
		\item \textbf{完备性质询 (Completeness)}：你的分类是否覆盖了所有情况，没有遗漏？
		
		\textit{反例}：你按“工作”、“生活”分类了所有文件，但一个“学习资料”的文件来了，它既不完全是工作，也不完全是生活。你的分类系统出现了漏洞（基底不完备），导致下次还得多花功夫，之前的分类也不确定是不是主要部分。
	\end{enumerate}
	
	\textbf{结论}：一个优秀的结构，必须满足 \lowfreq{MECE原则}（Mutually Exclusive, Collectively Exhaustive）。在数学上，这等价于寻找一组\textbf{正交且完备的基}。
	
	
	\newpage
	
	\subsection{分类即压缩：一种信息论视角}
	
	为什么建立结构（分类）能提升效率？因为它是一种极致的\lowfreq{信息压缩}。
	
	\begin{thinkbox}{“红色” vs “650纳米波长”：编码的层级}
		这两个词都可以指代同一种颜色，但它们的编码策略截然不同。
		
		\begin{center}
			\begin{tabular}{l|c|c}
				\toprule
				\textbf{特征} & \textbf{“红色” (粗分类)} & \textbf{“650nm” (精分类)} \\
				\midrule
				使用场景 & 日常生活 & 科学实验 \\
				出现频率 & \textbf{高频} & \textbf{低频} \\
				编码长度 & \textbf{短编码} (一个字) & \textbf{长编码} (多个字符+单位) \\
				误差容忍 & 容忍度高 & 极其敏感 \\
				\bottomrule
			\end{tabular}
		\end{center}
		
		\textbf{洞察}：
		\begin{itemize}
			\item \textbf{结构化就是选择合适的分类粒度}。对普通人说“红色”，信息传输极快，因为这是双方共有的、高频的短编码。
			\item 对科学家说“650nm”，是为了在要求极高精度的低频场景下，避免误差累积。
		\end{itemize}
		这统一了“结构化”和“变长编码”的底层逻辑：\textbf{都是为了在特定场景下，用最低的成本传递最有效的信息}。
	\end{thinkbox}
	
	我们已经确立了结构是根基，并用正交、完备的原则去审视它。接下来，我们将顺着这条逻辑链，深入探索信息传递的另一端——那个拥有着有限注意力、需要动机引导、并且依赖我们建立的这套高效编码系统来理解世界的“读者”。
	
	% = ... [上一部分代码结束] ...
	% 我们已经确立了结构是根基...深入探索...那个...“读者”。
	
	\newpage
	% ==============================================================================
	% Part 1: 读者端 - 机器的限制与欲望 (The Receiver)
	% ==============================================================================
	\section{ 读者端：机器的限制与欲望}
	
	我们已经搭建好了信息的“骨架”（结构），现在需要考虑信息要流向的终点——读者。分析读者，必须从两个正交的维度入手：他\textbf{能不能}读下去（生理限制），和他\textbf{想不想}读下去（心理动机）。
	
	\subsection{客观限制：认知能量预算 (Cognitive Energy Budget)}
	
	\begin{thinkbox}{为什么大脑会“累”？}
		大脑不是一个永动机。每一个需要深度思考的概念，每一次对模糊句子的解读，都在消耗宝贵的 ATP（三磷酸腺苷）。当能量消耗过快，保护机制就会启动，表现为：走神、困倦、放弃阅读。
		
		作者的任务，就是成为一个优秀的\lowfreq{能耗优化工程师}。
	\end{thinkbox}
	
	\begin{strategybox}{能量优化策略：变长编码的再应用}
		正如我们在开篇讨论的，面对有限的能量预算，最优策略就是变长编码。
		\begin{itemize}
			\item \textbf{为读者降噪}：对于高频、常见、符合直觉的信息，用最简短的笔墨带过。
			\item \textbf{强制聚焦}：对于低频、反直觉、高信息量的核心概念，比如\lowfreq{MECE原则}，必须使用加粗、加大字号等视觉手段，强制分配更多的认知能量到此处。
		\end{itemize}
		本文档中所有 \texttt{\textbackslash lowfreq} 的使用，都是这一原则的刻意为之。
	\end{strategybox}
	
	\subsection{主观动机：制造认知缺口 (Creating a Cognitive Gap)}
	
	\begin{thinkbox}{就算读者精力充沛，他为什么要选择读你的文章？}
		在一个信息过载的世界，注意力是一种主动的选择。仅仅降低阅读的能耗是不够的，我们还必须提供一个持续消耗能量的\textbf{理由}。如果读者没有阅读的意愿，再好的文章也是 0。
	\end{thinkbox}
	
	\begin{strategybox}{动机注入策略：主动提问}
		最高效的动机注入手段，不是给出答案，而是\textbf{提出一个好问题}。
		\begin{itemize}
			\item \textbf{制造缺口}：一个精准的问题，会在读者心中制造一个“知识缺口”。
			\item \textbf{驱动填补}：人类天生有填补认知缺口的欲望。为了消除这个“不确定性”带来的焦虑，他会主动读下去，寻找答案。
		\end{itemize}
		本文档中所有“思考链”模块的存在，其首要目的就是扮演这个“提问者”的角色，为你提供持续阅读的燃料。
	\end{strategybox}
	
	\subsection{通用桥梁：类比作为增量编码}
	
	当我们面对一个极其复杂的概念，逻辑链条太长，直接解释会同时耗尽读者的能量和动机。此时，我们需要一个“作弊”工具。
	
	\begin{thinkbox}{如何用最低成本解释一个新事物？}
		答案是：\textbf{别解释“新事物”，去解释它和“旧事物”的区别}。
		\lowfreq{“类比”}	在这里不是一种修辞手法，而是一种高效的\lowfreq{增量编码 (Incremental Encoding)} 算法。
		$$ \text{新知识} = \text{旧知识}_{\text{（读者脑中已有的短编码）}} + \Delta_{\text{（你需要传递的新信息）}} $$
		通过调用读者脑中已有的、高度压缩的“旧知识”作为基底，我们只需传递一个微小的增量 $\Delta$。这极大地降低了传输成本。
	\end{thinkbox}
	
	\subsection{自我质询：本章结构的正交与完备性}
	
	现在，让我们用开篇学到的原则，来审视本章自身的结构。
	
	\begin{strategybox}{MECE 自我审计}
		\begin{itemize}
			\item \textbf{正交性质询}：“认知能量”和“主观动机”是否独立？
			
			\textbf{回答}：是。前者是机器的\textbf{硬件算力}（能不能算），后者是操作系统的\textbf{任务指令}（要不要算）。一个顶配CPU（精力充沛）可以因为没有任务而待机（缺乏动机），一个老旧CPU（精力不足）也可以因为任务紧急而强行过载（动机强烈）。二者是正交维度。
			
			\item \textbf{完备性质询}：分析读者，只考虑这两点就够了吗？
			
			\textbf{回答}：对于写作这个宏观模型，已经足够。它覆盖了“意愿”和“能力”两个核心要素，构成了分析读者行为的完备基础。更细分的因素（如情绪、环境）可以被归纳到这两个大类之下。
		\end{itemize}
		
		通过这次质询，我们确保了本章的讨论是结构化、无冗余、无遗漏的。这再次印证了\textbf{结构优先}的原则。
	\end{strategybox}
	
	至此，我们完成了对“接收端”的分析。接下来，我们将转向信息的“发送端”，探讨一个同样重要的问题：在满足读者所有苛刻要求的同时，作者如何实现高效、可持续的内容生产？
	
	% = ... [上一部分代码结束] ...
	% 接下来，我们将转向信息的“发送端”，探讨...作者如何实现高效、可持续的内容生产？
	
	\newpage
	% ==============================================================================
	% Part 2: 作者端 - 生产与维护的工程学 (The Producer)
	% ==============================================================================
	\section{作者端：生产与维护的工程学}
	
	我们已经为读者铺设了一条认知上最高效的路径。现在，视角切换到我们自己——作者。如果为了服务读者而把自己累死，这个系统就是不可持续的。作者的效率同样至关重要。
	
	\subsection{核心痛点：修改成本的不对称性}
	
	\begin{thinkbox}{为什么改一个错别字的成本，和改一个核心论点的成本，天差地别？}
		写作不是一次成型的过程，而是一个不断迭代、修改的过程。然而，不同修改行为带来的痛苦是极不均衡的。
		\begin{itemize}
			\item \textbf{修改细节 (Detail)}：改一个形容词，修正一个数据。影响范围极小，成本可忽略。
			\item \textbf{修改框架 (Framework)}：调整章节顺序，更换核心论点。这会引发连锁反应，导致底下所有细节全部作废重写。成本是灾难性的。
		\end{itemize}
		
		这就是写作中的\lowfreq{修改不等式 (Modification Inequality)}：
		$$ \text{Cost}(\Delta \text{Framework}) \gg \text{Cost}(\Delta \text{Detail}) $$
	\end{thinkbox}
	
	这个不等式告诉我们，为了避免返工，我们必须采用一种能将“框架”和“细节”隔离的生产方式。这个方式，就是软件工程的核心思想——模块化。
	
	\subsection{解决方案：万物皆可模块化 (Modularity)}
	
	\lowfreq{模块化}是将一个复杂的系统，拆分成多个高内聚、低耦合的独立单元的艺术。
	
	\subsubsection{内部模块化：黑盒与接口 (Black Box \& Interface)}
	
	我们要把文章的每一个章节，都看作是一个独立的“函数”或“模块”。
	
	\begin{strategybox}{模块化写作流程}
		\begin{enumerate}
			\item \textbf{定义接口 (Interface)}：明确本章节的输入（依赖于前文的哪些结论）和输出（要向后文传递什么结论）。接口一旦定义，就轻易不要改动。
			\item \textbf{黑盒实现 (Implementation)}：在接口不变的前提下，章节内部的论证方式、使用的例子、具体的措辞，都可以随意修改，而不会影响到任何其他章节。
		\end{enumerate}
		这种设计，将高成本的“框架修改”锁定在了少数几个“接口”定义上，而将绝大部分的修改工作，都变成了低成本的“内部实现”调整。
	\end{strategybox}
	
	\subsubsection{外部模块化：可扩展性插槽 (Extensibility)}
	
	任何一篇文章都不可能在第一版就达到“完备”。总有我们暂时没想到，或者需要未来补充的内容。
	
	\begin{strategybox}{预留扩展接口}
		一个健壮的系统，必须为未来留白。与其假装自己无所不知，不如预先设计好“扩展插槽”。
		
		例如，在本文档的结构设计中，我们可以预留一个关于“协作写作中的版本控制”的插槽。现在不必写，但框架已经为它留好了位置。这保证了结构的长期完备性。
	\end{strategybox}
	
	\subsection{自我质询：本章结构再次审计}
	我们再次用 MECE 原则审视自身。
	
	\begin{strategybox}{MECE 自我审计}
		\begin{itemize}
			\item \textbf{正交性质询}：“修改不等式”（问题）和“模块化”（解决方案）是否独立？
			
			\textbf{回答}：是。前者是对一个客观现象的描述，后者是针对该现象提出的一套工程方法论。它们构成了标准的“问题-解决”正交对。
			
			\item \textbf{完备性质询}：作者端的效率问题，只靠模块化就够了吗？
			
			\textbf{回答}：在“结构”这个宏观层面，模块化是核心。其他效率工具（如自动化排版、文献管理）都是在这个模块化思想下的具体实现，可以被归纳到这个框架内。因此，对于本理论体系而言，它是完备的。
		\end{itemize}
	\end{strategybox}
	
	至此，我们构建了从“顶层结构设计”，到“读者端接收优化”，再到“作者端生产效率”的完整闭环。所有这些工程化原则，都建立在坚实的逻辑和数学基础之上。
	
	所有这些原则背后的硬核数学证明，都统一收纳于最终的附录之中，那里是本理论体系的“发动机室”。
	
	% = ... [上一部分代码结束] ...
	% 那里是本理论体系的“发动机室”。
	
	\newpage
	% ==============================================================================
	% 附录：深度数学推导 (The Engine Room)
	% ==============================================================================
	\appendix
	\part*{附录：支撑理论的数学基石}
	\addcontentsline{toc}{part}{附录：支撑理论的数学基石}
	
	这里存放着所有“枯燥”但“绝对正确”的数学推导。我们将它们从正文中剥离，是为了保证正文的阅读流畅性，同时为所有工程原则提供终极的、无可辩驳的严谨性证明。
	
	% --- Math 1: 变长编码 ---
	\section{A. 变长编码的拉格朗日推导} \label{math:vlc}
	\mathlink{math:vlc}{正文链接：为何低频必须长编码？}
	
	\begin{mathbox}{1. 问题动机与写作场景映射}
		\textbf{问题}：在写作中，我们凭什么说“重点内容要浓墨重彩”？这背后是否有最优化的数学依据？
		
		\textbf{映射}：我们将“读者的注意力/认知能量”建模为一个需要被最小化的成本函数。一篇文章由不同“信息量”的词汇组成，“篇幅/视觉权重”就是我们要分配的资源。目标是，在保证信息完整传达的前提下，让读者总的阅读成本最低。
	\end{mathbox}
	
	\begin{mathbox}{2. 建模过程 (Modeling)}
		\begin{itemize}
			\item 设词汇表中有 $M$ 个词 $w_1, \dots, w_M$。
			\item $p_i$：第 $i$ 个词出现的概率（频率）。\textbf{高频词}如“的”，$p_i$ 大；\textbf{低频核心概念}如“正交化”，$p_i$ 小。
			\item $l_i$：分配给第 $i$ 个词的“编码长度”（视觉权重或篇幅）。这是我们的\textbf{决策变量}。
		\end{itemize}
		
		\textbf{目标函数 (Objective)}：最小化平均阅读成本（期望长度）。
		$$ \min \bar{L} = \sum_{i=1}^{M} p_i l_i $$
		
		\textbf{约束条件 (Constraint)}：信息必须能被无歧义地解码。这由 Kraft 不等式保证。
		$$ \sum_{i=1}^{M} 2^{-l_i} \le 1 $$
	\end{mathbox}
	
	\newpage
	
	\begin{mathbox}{3. 数学公式每一步的动机解释 (Step-by-Step Derivation)}
		\textbf{动机}：这是一个典型的带约束优化问题，使用\lowfreq{拉格朗日乘子法}是标准解法。
		
		\textbf{Step 1}: 构造拉格朗日函数 $\mathcal{L}$，引入乘子 $\lambda$ 将约束纳入目标函数。
		$$ \mathcal{L}(l_i, \lambda) = \sum p_i l_i + \lambda \left( \sum 2^{-l_i} - 1 \right) $$
		
		\textbf{Step 2}: 寻找极值点。我们对每个决策变量 $l_i$ 求偏导，并令其为 0。
		$$ \frac{\partial \mathcal{L}}{\partial l_i} = \frac{\partial}{\partial l_i} (p_i l_i) + \lambda \frac{\partial}{\partial l_i} (2^{-l_i}) = 0 $$
		$$ p_i + \lambda (2^{-l_i} \ln 2 \cdot (-1)) = 0 $$
		
		\textbf{Step 3}: 整理方程，解出 $l_i$。
		$$ p_i = \lambda \ln 2 \cdot 2^{-l_i} $$
		$$ 2^{-l_i} = \frac{p_i}{\lambda \ln 2} $$
		两边取 $\log_2$，得到：
		$$ -l_i = \log_2(p_i) - \log_2(\lambda \ln 2) $$
		$$ l_i = -\log_2(p_i) + C \quad (\text{其中 } C = \log_2(\lambda \ln 2) \text{ 是一个常数}) $$
	\end{mathbox}
	
	\begin{mathbox}{4. 结论与写作中的映射}
		\textbf{数学结论}：最优的编码长度 $l_i$ 与信息源概率 $p_i$ 的对数成反比。这正是信息论中\lowfreq{香农信息量}的定义。
		$$ l_i \approx H(p_i) = -\log_2 p_i $$
		
		\textbf{写作映射}：
		\begin{itemize}
			\item 对于高频词（如“的”，$p_i \to 1$），其信息量 $H(p_i) \to 0$。我们应该用最短的篇幅（即标准字体）处理。
			\item 对于低频核心概念（如“拉格朗日”，$p_i \to 0$），其信息量 $H(p_i) \to \infty$。我们必须用最长的篇幅（加大、加粗、反复解释）来匹配其巨大的信息量。
		\end{itemize}
		这为我们在正文中使用的 \texttt{\textbackslash lowfreq} 宏提供了坚实的数学基础。
	\end{mathbox}
	
	% --- 为了视觉分隔，可以添加一个分页 ---
	\newpage
	
	% --- Math 2: 正交化 ---
	\section{B. 逻辑去重的 Gram-Schmidt 过程} \label{math:ortho}
	\mathlink{math:ortho}{正文链接：如何从数学上剔除重复逻辑？}
	
	\begin{mathbox}{1. 问题动机与写作场景映射}
		\textbf{问题}：当我们进行头脑风暴时，会产生大量想法。如何系统性地整理这些想法，去除重复的“废话”，并保留真正独立、有价值的论点？
		
		\textbf{映射}：我们将每一个“想法”或“论点”建模为高维空间中的一个\textbf{向量}。如果两个想法本质上在说一件事，那么这两个向量就是\textbf{线性相关}的。我们的目标，就是从这组杂乱的向量中，提取出一组\lowfreq{正交基}。
	\end{mathbox}
	
	\begin{mathbox}{2. 建模与推导 (Gram-Schmidt Process)}
		\textbf{输入}：一组线性相关的论点向量 $\{v_1, v_2, \dots, v_n\}$。
		\textbf{输出}：一组正交的论点向量 $\{u_1, u_2, \dots, u_n\}$。
		
		\textbf{Step 1}: 选定第一个论点作为基准。
		$$ u_1 = v_1 $$
		
		\textbf{Step 2}: 处理第二个论点 $v_2$。我们首先计算 $v_2$ 在 $u_1$ 方向上的\textbf{投影}，这个投影就是 $v_2$ 中与 $u_1$ 重复的“废话”部分。
		$$ \text{proj}_{u_1}(v_2) = \frac{\langle v_2, u_1 \rangle}{\langle u_1, u_1 \rangle} u_1 $$
		然后，从 $v_2$ 中减去这个“废话”部分，剩下的就是纯粹的、与 $u_1$ 正交的新信息。
		$$ u_2 = v_2 - \text{proj}_{u_1}(v_2) $$
		
		\textbf{Step k (通用步骤)}: 对于第 $k$ 个论点 $v_k$，我们减去它在所有已经正交化的基 $\{u_1, \dots, u_{k-1}\}$ 上的投影。
		$$ u_k = v_k - \sum_{j=1}^{k-1} \text{proj}_{u_j}(v_k) $$
	\end{mathbox}
	
	\begin{mathbox}{3. 结论与写作中的映射}
		\textbf{数学结论}：Gram-Schmidt 过程提供了一个算法化的流程，可以将任何一组线性无关的向量转化为一组标准正交基。
		
		\textbf{写作映射}：这为我们提供了审视文章结构的“手术刀”。
		\begin{itemize}
			\item 在列提纲时，每写下一个新的标题，都应该在头脑中进行一次“投影”操作：这个新论点，有多少成分是前面已经说过的？
			\item 如果投影分量很大，说明这个论点是冗余的，应该合并或删除。
			\item 只有当一个论点减去所有“投影”后，剩下的“残差”向量仍然足够大，它才有资格成为一个独立的章节。
		\end{itemize}
		这正是 \lowfreq{MECE原则} 中“Mutually Exclusive”（相互独立）的数学化身。
	\end{mathbox}
	
	
	
\end{document} ==========================================================